% ===========================================================
% -----------------------------------------------------------
% Resumo, Abastract , Résumé  e Resumen
% -----------------------------------------------------------
% ===========================================================

% -----------------------------------------------------------
% 	Resumo em português
% -----------------------------------------------------------
\setlength{\absparsep}{18pt} % ajusta o espaçamento dos parágrafos do resumo
\begin{resumo}
	O diabetes é uma doença crônica grave e prevalente, especialmente o Tipo 2, associada ao estilo de vida e que representa um desafio crescente para a saúde pública global. Neste contexto, as redes neurais artificiais surgem como uma ferramenta promissora para a detecção e prevenção, analisando grandes volumes de dados clínicos e comportamentais para identificar padrões complexos e classificar com precisão o risco individual. A integração desses modelos em sistemas de apoio à decisão médica potencia diagnósticos mais rápidos e o desenvolvimento de estratégias de intervenção precoce, visando reduzir a incidência e o impacto da doença.
	
	\textbf{Palavras-chave}: Redes neurais, diabetes, predição, inteligência artificial, aprendizado de máquina.
\end{resumo}

% -----------------------------------------------------------
% 	Resumo em inglês
% -----------------------------------------------------------
\begin{resumo}[Abstract]
	\begin{otherlanguage*}{english}
		Diabetes is a serious and prevalent chronic disease, particularly Type 2, which is often linked to lifestyle factors and represents a growing challenge for global public health. In this context, artificial neural networks are emerging as a promising tool for detection and prevention, capable of analyzing large volumes of clinical and behavioral data to identify complex patterns and accurately classify an individual's risk. Integrating these models into clinical decision support systems enhances the potential for faster diagnoses and the development of early intervention strategies, aiming to reduce the disease's incidence and impact.
		
		\vspace{\onelineskip}
		
		\noindent 
		\textbf{Keywords}: Neural networks, diabetes, prediction, artificial intelligence, machine learning.

	\end{otherlanguage*}
\end{resumo}
